\vspace{0.5em}\noindent\textbf{Hệ gợi ý sản phẩm thương mại điện tử trên
AWS}\par

\vspace{0.5em}\noindent\textbf{Website Thương mại điện tử Tích hợp Gợi ý Cá nhân hóa và Tìm
kiếm bằng Hình ảnh trên
AWS}\par

\vspace{0.5em}\noindent\textbf{1. Tóm tắt điều
hành}\par

Dự án \textbf{Hệ gợi ý sản phẩm thương mại điện tử trên AWS} được thực
hiện trong khuôn khổ đồ án môn \textbf{Machine Learning trên AWS}. Mục
tiêu của dự án là xây dựng và triển khai một trang web thương mại điện
tử hoàn chỉnh, tích hợp hệ thống gợi ý sản phẩm cá nhân hóa và tính năng
tìm kiếm bằng hình ảnh. Nền tảng tận dụng các dịch vụ AWS Serverless
(API Gateway, AWS Lambda, DynamoDB, S3, CloudFront) kết hợp với Amazon
Personalize và mô hình CLIP để mang lại trải nghiệm mua sắm tối ưu và
hiện đại.

\vspace{0.5em}\noindent\textbf{2. Tuyên bố vấn
đề}\par

\emph{Vấn đề hiện tại}\\
Các website thương mại điện tử truyền thống thường thiếu khả năng cá
nhân hóa trải nghiệm theo thời gian thực, dẫn đến tỷ lệ chuyển đổi và độ
tương tác của người dùng chưa cao. Bên cạnh đó, việc tìm kiếm sản phẩm
bằng từ khóa văn bản bị hạn chế khi người dùng chỉ có hình ảnh sản phẩm.
Chi phí vận hành hạ tầng máy chủ truyền thống cũng là một rào cản lớn
khi hệ thống cần mở rộng quy mô.

\emph{Giải pháp}\\
Hệ thống giải quyết các vấn đề trên bằng cách triển khai giải pháp
e-commerce hoàn chỉnh trên hạ tầng Serverless của AWS:

\begin{itemize}
\tightlist
\item
  \textbf{Gợi ý cá nhân hóa:} Tận dụng Amazon Personalize (recipe
  \texttt{aws-user-personalization}) để hiển thị khối ``Gợi ý riêng cho
  bạn'' ngay tại trang chủ.
\item
  \textbf{Tìm kiếm bằng hình ảnh:} Sử dụng mô hình CLIP
  (\texttt{@xenova/transformers}) kết hợp đo độ tương đồng Cosine
  (Cosine Similarity).
\item
  \textbf{Nghiệp vụ Mua sắm Toàn diện:} Duyệt/tìm kiếm/lọc sản phẩm,
  quản lý tài khoản/phiên làm việc, giỏ hàng, mã giảm giá, đặt hàng (COD
  / Chuyển khoản) và đánh giá sản phẩm.
\item
  \textbf{Kiến trúc Serverless:} Giao diện React (Vite) lưu trữ trên S3
  và phân phối qua CloudFront, API Node.js 20 chạy trên AWS Lambda qua
  API Gateway, cơ sở dữ liệu DynamoDB (8 bảng) hoạt động theo chế độ
  On-demand.
\end{itemize}

\emph{Lợi ích và hoàn vốn đầu tư (ROI)}

\begin{itemize}
\tightlist
\item
  \textbf{Tăng tỷ lệ chuyển đổi:} Gợi ý sản phẩm chính xác dựa trên hành
  vi giúp tăng giá trị đơn hàng và độ hài lòng của khách hàng.
\item
  \textbf{Tối ưu chi phí hạ tầng:} Mô hình Serverless giúp chỉ trả tiền
  theo lượng truy cập thực tế, loại bỏ chi phí duy trì máy chủ rảnh rỗi.
\item
  \textbf{Khả năng mở rộng cao:} DynamoDB On-demand và Lambda tự động co
  giãn theo tải mà không cần quản trị hạ tầng thủ công.
\end{itemize}

\vspace{0.5em}\noindent\textbf{3. Kiến trúc giải
pháp}\par

Nền tảng sử dụng kiến trúc Serverless trên AWS để xử lý toàn bộ luồng
nghiệp vụ và gợi ý AI:

\begin{figure}
\centering
\pandocbounded{\includegraphics[keepaspectratio,alt={Mô tả hình ảnh}]{5-Workshop/architecture.jpg}}
\caption{Mô tả hình ảnh}
\end{figure}

\emph{Dịch vụ AWS \& Công nghệ sử dụng}

\begin{itemize}
\tightlist
\item
  \textbf{AWS CloudFront \& Amazon S3:} Lưu trữ giao diện tĩnh React 18
  (Vite, Redux Toolkit, Tailwind) và phân phối tại biên.
\item
  \textbf{Amazon API Gateway:} HTTP API quản lý 22 routes chuyển tiếp
  tới AWS Lambda.
\item
  \textbf{AWS Lambda:} Hàm \texttt{fcj-api} (Node.js 20) xử lý toàn bộ
  logic API và kết nối cơ sở dữ liệu.
\item
  \textbf{Amazon DynamoDB:} 8 bảng dữ liệu (\texttt{Products},
  \texttt{Categories}, \texttt{Users}, \texttt{Sessions},
  \texttt{Carts}, \texttt{Vouchers}, \texttt{Reviews}, \texttt{Orders})
  ở chế độ On-demand.
\item
  \textbf{Amazon Personalize:} Sử dụng recipe
  \texttt{aws-user-personalization} cung cấp danh sách gợi ý theo
  Campaign ARN.
\item
  \textbf{Tìm kiếm bằng hình ảnh:} CLIP (\texttt{@xenova/transformers})
  + Cosine Similarity (chạy ở máy local).
\end{itemize}

\emph{Thiết kế thành phần}

\begin{itemize}
\tightlist
\item
  \textbf{Frontend:} Code build được tải lên S3 bucket riêng tư, chỉ
  CloudFront được phép truy cập qua OAC.
\item
  \textbf{Backend API:} Lambda tiếp nhận 22 endpoint nghiệp vụ, băm mật
  khẩu người dùng bằng \texttt{scrypt} và so sánh an toàn bằng
  \texttt{timingSafeEqual}.
\item
  \textbf{Hệ gợi ý:} Lambda gọi API tới Personalize Campaign để lấy danh
  sách gợi ý riêng cho từng người dùng đã đăng nhập.
\item
  \textbf{Tìm kiếm hình ảnh:} Tạo vector từ ảnh sản phẩm lưu tại
  \texttt{vector-cache.json}, so sánh vector đường dẫn ảnh đầu vào để
  trả về sản phẩm tương đồng.
\end{itemize}

\vspace{0.5em}\noindent\textbf{4. Triển khai kỹ
thuật}\par

\emph{Các giai đoạn triển khai}

\begin{enumerate}
\def\labelenumi{\arabic{enumi}.}
\tightlist
\item
  \textbf{Khởi tạo Hạ tầng \& Database:} Thiết kế schema cho 8 bảng
  DynamoDB, cấu hình IAM Role, S3, CloudFront và API Gateway.
\item
  \textbf{Phát triển Backend \& Frontend:} Viết 22 route API trong
  \texttt{backend/index.mjs}, dựng giao diện React và kết nối Redux
  Toolkit với API thật.
\item
  \textbf{Tích hợp ML \& AI:} Xây dựng pipeline dữ liệu
  (\texttt{ml\_recommendation\_dataset.csv}), huấn luyện mô hình Amazon
  Personalize và cài đặt CLIP visual search.
\item
  \textbf{Kiểm thử \& Triển khai:} Viết kịch bản kiểm thử tự động
  Playwright (luồng đặt hàng), deploy Lambda \& CloudFront, đánh giá chỉ
  số mô hình.
\end{enumerate}

\emph{Yêu cầu kỹ thuật}

\begin{itemize}
\tightlist
\item
  \textbf{Môi trường:} Node.js 20 trở lên, AWS CLI / IAM User được cấp
  đúng quyền trên 8 bảng DynamoDB và Personalize Campaign.
\item
  \textbf{Bảo mật:} Không commit file \texttt{.env} chứa Access Key lên
  Git; phân quyền IAM theo nguyên tắc tối thiểu; băm mật khẩu người dùng
  bằng \texttt{scrypt}.
\end{itemize}

\vspace{0.5em}\noindent\textbf{5. Lộ trình \& Mốc triển
khai}\par

\begin{itemize}
\tightlist
\item
  \textbf{Giai đoạn 1: Lập kế hoạch \& Xây dựng dữ liệu}

  \begin{itemize}
  \tightlist
  \item
    Thiết kế kiến trúc AWS và bảng cơ sở dữ liệu.
  \item
    Sinh bộ dữ liệu tương tác mô phỏng hành vi mua sắm (phễu chuyển đổi,
    quy luật power-law).
  \end{itemize}
\item
  \textbf{Giai đoạn 2: Phát triển Core E-Commerce}

  \begin{itemize}
  \tightlist
  \item
    Hoàn thiện 22 route API trên Lambda.
  \item
    Xây dựng giao diện frontend (Trang chủ, Chi tiết sản phẩm, Giỏ hàng,
    Thanh toán, Tài khoản).
  \end{itemize}
\item
  \textbf{Giai đoạn 3: Tích hợp Machine Learning}

  \begin{itemize}
  \tightlist
  \item
    Huấn luyện và tạo Campaign trên Amazon Personalize.
  \item
    Tích hợp tính năng tìm kiếm bằng hình ảnh CLIP.
  \end{itemize}
\item
  \textbf{Giai đoạn 4: Kiểm thử \& Triển khai}

  \begin{itemize}
  \tightlist
  \item
    Viết kịch bản Playwright test tự động luồng checkout.
  \item
    Deploy backend lên Lambda và frontend lên S3/CloudFront.
  \end{itemize}
\end{itemize}

\vspace{0.5em}\noindent\textbf{6. Ước tính ngân sách \& Chỉ số Mô
hình}\par

\textbf{Ước tính chi phí AWS hàng tháng (Quy mô lưu lượng thấp / Đồ án
môn học)}

{\def\LTcaptype{none} % do not increment counter
\begingroup
\small
\setlength{\tabcolsep}{3pt}
\renewcommand{\arraystretch}{1.15}
\begin{longtable}[]{@{}
  >{\raggedright\arraybackslash}p{(\linewidth - 4\tabcolsep) * \real{0.3697}}
  >{\raggedright\arraybackslash}p{(\linewidth - 4\tabcolsep) * \real{0.3950}}
  >{\raggedright\arraybackslash}p{(\linewidth - 4\tabcolsep) * \real{0.2353}}@{}}
\toprule\noalign{}
\begin{minipage}[b]{\linewidth}\raggedright
Dịch vụ AWS
\end{minipage} & \begin{minipage}[b]{\linewidth}\raggedright
Giả định mức sử dụng
\end{minipage} & \begin{minipage}[b]{\linewidth}\raggedright
Chi phí ước tính (USD/tháng)
\end{minipage} \\
\midrule\noalign{}
\endhead
\bottomrule\noalign{}
\endlastfoot
\textbf{AWS Lambda} (\texttt{fcj-api}) & \textasciitilde1 triệu lượt
gọi, 512MB, \textasciitilde200ms trung bình & \$0.20 -- \$1.00 \\
\textbf{API Gateway (HTTP API)} & \textasciitilde1 triệu request &
\textasciitilde\$1.00 \\
\textbf{Amazon DynamoDB} (8 bảng, On-demand) & \textasciitilde1 triệu
request unit đọc/ghi & \$5 -- \$10 \\
\textbf{Amazon S3} (lưu trữ frontend) & \textasciitilde5 GB dung lượng +
request PUT/GET & \textasciitilde\$0.15 \\
\textbf{Amazon CloudFront} & \textasciitilde50 GB dữ liệu truyền tải ra
& \textasciitilde\$4.25 \\
\textbf{Amazon Personalize (Campaign)} & 1 campaign, TPS tối thiểu, chạy
24/7 (\textasciitilde730 giờ) & \textbf{\$150 -- \$220} \\
\textbf{Amazon Personalize (Huấn luyện Solution)} & Huấn luyện lại định
kỳ (\textasciitilde2--4 giờ/lần) & \textasciitilde\$0.50 -- \$1.00 mỗi
lần \\
\textbf{CloudWatch Logs \& Monitoring} & Gói ghi log cơ bản &
\textasciitilde\$1.00 \\
\textbf{Tổng cộng (ổn định, hàng tháng)} & & \textbf{≈ \$160 -- \$240 /
tháng} \\
\end{longtable}
\endgroup
}

\begin{quote}
\textbf{Lưu ý:} Amazon Personalize Campaign chạy liên tục là yếu tố chi
phí lớn nhất, do tính phí theo TPS tối thiểu mỗi giờ bất kể lưu lượng
truy cập thực tế. Đối với đồ án môn học/demo, có thể giảm đáng kể chi
phí này bằng cách xóa Campaign khi không sử dụng để demo, hoặc chỉ tạo
lại trong giai đoạn đánh giá/kiểm thử. Chi phí Lambda, API Gateway,
DynamoDB, S3 và CloudFront ở mức tối thiểu tại quy mô lưu lượng này nhờ
mô hình tính phí Serverless theo mức sử dụng thực tế.
\end{quote}

\emph{Đánh giá Bộ dữ liệu Huấn luyện (v1 vs v2)} Bộ dữ liệu tương tác
được sinh mô phỏng theo hành vi mua sắm thực tế mang lại kết quả cải
thiện vượt trội:

{\def\LTcaptype{none} % do not increment counter
\begingroup
\small
\setlength{\tabcolsep}{3pt}
\renewcommand{\arraystretch}{1.15}
\begin{longtable}[]{@{}lll@{}}
\toprule\noalign{}
Chỉ số & Dữ liệu v1 & Dữ liệu v2 \\
\midrule\noalign{}
\endhead
\bottomrule\noalign{}
\endlastfoot
\textbf{Precision@5} & 0.0889 & 0.4348 \\
\textbf{NDCG@10} & 0.1799 & 0.6512 \\
\textbf{MRR@25} & 0.1216 & 0.7130 \\
\textbf{Coverage} & 0.8218 & 0.9505 \\
\end{longtable}
\endgroup
}

\emph{Kết luận:} Chất lượng dữ liệu đầu vào đóng vai trò quyết định đến
hiệu năng mô hình hơn việc thay đổi thuật toán.

\vspace{0.5em}\noindent\textbf{7. Đánh giá rủi ro}\par

\emph{Ma trận rủi ro}

\begin{itemize}
\tightlist
\item
  \textbf{Vòng lặp dữ liệu chưa khép kín:} Ảnh hưởng cao, xác suất trung
  bình (Tương tác thực của người dùng chưa được ghi ngược lại tự động để
  retraining định kỳ).
\item
  \textbf{Giới hạn kích thước Lambda:} Ảnh hưởng trung bình, xác suất
  cao (Model CLIP nặng vài trăm MB không thể deploy trực tiếp lên
  Lambda, phải chạy local).
\item
  \textbf{Cache CloudFront bị cũ:} Ảnh hưởng thấp, xác suất cao (Cần
  invalidate cache \texttt{/*} sau mỗi lần upload frontend mới lên S3).
\item
  \textbf{Truy vấn dữ liệu phức tạp (Database Bottleneck):} Ảnh hưởng
  trung bình, xác suất cao (DynamoDB không tối ưu cho việc tìm kiếm văn
  bản tự do (full-text search) và lọc sản phẩm theo nhiều tiêu chí).
\end{itemize}

\emph{Chiến lược giảm thiểu}

\begin{itemize}
\tightlist
\item
  \textbf{Dữ liệu:} Nghiên cứu tích hợp AWS EventBridge / Kinesis để ghi
  log tương tác phục vụ retrain tự động.
\item
  \textbf{Model nặng:} Hướng tới sử dụng AWS Fargate hoặc Amazon
  SageMaker Endpoint cho mô hình tìm kiếm hình ảnh.
\item
  \textbf{Deploy:} Tự động hóa lệnh invalidate CloudFront trong script
  deploy.
\item
  \textbf{Truy vấn:} Kết hợp Amazon OpenSearch để xử lý các truy vấn tìm
  kiếm và lọc sản phẩm phức tạp.
\end{itemize}

\vspace{0.5em}\noindent\textbf{8. Kết quả kỳ
vọng}\par

\emph{Cải tiến kỹ thuật:} Hệ thống thương mại điện tử hoàn chỉnh chạy
trên Serverless AWS, tích hợp kiểm thử tự động Playwright, phản hồi gợi
ý cá nhân hóa dưới 1 giây. \emph{Giá trị dài hạn:} Làm mô hình tham
chiếu chuẩn cho việc triển khai các ứng dụng Machine Learning thực tế
lên nền tảng đám mây AWS.

\vspace{0.5em}\noindent\textbf{9. Một số hướng cải
thiện}\par

Dựa trên thực tế vận hành và yêu cầu mở rộng, một số dịch vụ AWS có thể
được thêm vào hệ thống nhằm tối ưu hóa bảo mật, hiệu năng và khả năng mở
rộng:

\begin{itemize}
\tightlist
\item
  \textbf{Bảo mật nhiều lớp:} Bổ sung AWS WAF trước CloudFront/API
  Gateway để chống Botnet/DDoS. Chuyển đổi quản lý người dùng sang
  Amazon Cognito.
\item
  \textbf{Microservices:} Tách nhỏ Lambda hiện tại thành các dịch vụ độc
  lập (User, Product, Order) để tránh rủi ro hệ thống và cải thiện thời
  gian khởi động (Cold Start).
\item
  \textbf{Tối ưu Truy vấn \& Tìm kiếm:} Tích hợp Amazon OpenSearch cùng
  với DynamoDB để đáp ứng các truy vấn lọc sản phẩm và text-search phức
  tạp. Đưa mô hình CLIP lên Amazon ECS hoặc SageMaker.
\item
  \textbf{Real-time Feedback Loop:} Sử dụng Kinesis Data Streams kết hợp
  với Amazon Personalize Event Tracker để cập nhật dữ liệu hành vi người
  dùng theo thời gian thực.
\end{itemize}
